%% =================================================================
%% Wei, Mei, Zhao, Xiao, Sun, Zhang, Guo.
%% "Nondestructively identifying the mechanical behavior of soft
%%  tissues using surface deformation with an explicit inverse
%%  approach."
%% Appl. Math. Modelling 134 (2024) 126--147.
%%
%% Reconstruction of page 12 (printed p. 137) as study notes.
%% Generated by: reproduce-page-basic (first held-out validation page)
%%
%% Structure: mirrors pages 9, 10 — two full-page result figures
%% (Fig. 12, Fig. 13) + one short prose paragraph.
%% Prose: extracted verbatim from the PDF text layer via pdftotext.
%% Math:  none on this page.
%% Figures 12 & 13: placeholder boxes with captions + visual descriptions.
%% =================================================================

\documentclass[11pt]{article}

\usepackage[margin=1in]{geometry}
\usepackage{amsmath, amssymb}
\usepackage{mathrsfs}
\usepackage{bm}
\usepackage{xcolor}
\usepackage{fancyhdr}

\pagestyle{fancy}
\fancyhf{}
\lhead{\small Y.~Mei et al.}
\rhead{\small Applied Mathematical Modelling 134 (2024) 126--147}
\cfoot{\small 137 $\to$ \arabic{page}}
\renewcommand{\headrulewidth}{0.4pt}

\setcounter{page}{1}

\begin{document}

% ---------- Figure 12 placeholder ----------
\noindent
\begin{center}
\fbox{\parbox{0.92\textwidth}{\centering\vspace{1em}
\textbf{[Figure 12 --- placeholder]}\\[0.8em]
Six-panel side-by-side comparison of the \textbf{Explicit Inverse
Method} (left column) vs.\ \textbf{Implicit Inverse Method} (right
column) on \textbf{Sample 2} (the two-layer target: background
$\mu^0 \approx 3$\,MPa, top layer $\mu^1 \approx 10$\,MPa),
with \textbf{1\,\% noise} added to the measured surface
displacements, using 3, 6, and 9 displacement datasets.\\[0.4em]
\textbf{Explicit (left column)}: bilayer structure recovered with
a relatively clean (though slightly wavy) interface. Colorbar
ranges:
(a) 3 fields, SM $\in [1.00, 11.33]$\,MPa;
(b) 6 fields, SM $\in [1.00, 11.43]$\,MPa;
(c) 9 fields, SM $\in [1.00, 11.34]$\,MPa.\\[0.3em]
\textbf{Implicit (right column)}: bilayer structure visible but
with significant artifacts, especially in panel (d) where a large
green/blue patch intrudes into the red top layer from the left
edge. Colorbar ranges:
(d) 3 fields, SM $\in [1.00, 14.60]$\,MPa --- a \textbf{46\,\%
overshoot} above the nominal 10\,MPa top layer, matching the
paper's claim that the implicit method's average relative error
``could exceed 40\,\%'' under 1\,\% noise;
(e) 6 fields, SM $\in [1.00, 10.60]$\,MPa;
(f) 9 fields, SM $\in [1.00, 10.71]$\,MPa.\\[0.3em]
\textbf{Takeaway}: the noise amplifies the implicit method's
difficulty with interface recovery dramatically in the 3-field
case; 6 and 9 fields bring it back toward the nominal value but
with visibly rougher interfaces than the explicit method.
\vspace{1em}}}
\end{center}
\vspace{0.3em}
\begin{center}
\textbf{Fig. 12.} The reconstructed results with varying 3, 6 and
9 surface displacement datasets corresponding to Sample 2 (1\,\%
noise is introduced into the measured data, Unit: MPa).
\end{center}

\vspace{0.8em}

% ---------- Figure 13 placeholder ----------
\noindent
\begin{center}
\fbox{\parbox{0.92\textwidth}{\centering\vspace{1em}
\textbf{[Figure 13 --- placeholder]}\\[0.8em]
Six-panel side-by-side comparison of the \textbf{Explicit Inverse
Method} (left column) vs.\ \textbf{Implicit Inverse Method} (right
column) on \textbf{Sample 3} (the three-layer target: blue bottom
layer $\mu^0 \approx 3$\,MPa, green middle layer
$\mu^1 \approx 6$\,MPa, red top layer $\mu^2 \approx 10$\,MPa),
with \textbf{1\,\% noise} added to the measured surface
displacements, using 3, 6, and 9 displacement datasets.\\[0.4em]
\textbf{Explicit (left column)}: three-layer structure recovered
with wavy but distinct interfaces between the layers. Colorbar
ranges:
(a) 3 fields, SM $\in [1.01, 10.63]$\,MPa;
(b) 6 fields, SM $\in [1.00, 10.25]$\,MPa;
(c) 9 fields, SM $\in [1.00, 10.38]$\,MPa.\\[0.3em]
\textbf{Implicit (right column)}: three-layer structure noticeably
noisier, with fuzzier interfaces and small artifact patches.
Colorbar ranges:
(d) 3 fields, SM $\in [1.00, 11.78]$\,MPa;
(e) 6 fields, SM $\in [1.00, 11.58]$\,MPa;
(f) 9 fields, SM $\in [1.00, 9.96]$\,MPa --- notably, the 9-field
implicit result \textbf{undershoots} the nominal 10\,MPa target,
unusual relative to most other panels in the paper which tend to
overshoot.\\[0.3em]
\textbf{Takeaway}: as with Sample 2 under noise, the explicit
method produces cleaner interfaces; but the quantitative
colorbar-max values on Sample 3 are closer between explicit and
implicit than on Sample 2 because the three-layer structure
constrains both methods more tightly than the two-layer structure.
\vspace{1em}}}
\end{center}
\vspace{0.3em}
\begin{center}
\textbf{Fig. 13.} The reconstructed results with varying 3, 6 and
9 surface displacement datasets corresponding to Sample 3 (1\,\%
noise is introduced into the measured data, Unit: MPa).
\end{center}

\vspace{1em}

% ---------- Prose paragraph (from pdftotext, soft-hyphen cleaned) ----------
\noindent
When 1\,\% noise was introduced, the explicit inverse method
continued to perform admirably, producing high-quality
reconstructed results (see Fig.~20). However, the performance of
the implicit inverse method continued to be inferior to that of
the explicit inverse method. The average relative error for the
implicit inverse method could exceed 40\,\%, indicating substantial
inaccuracies in the reconstructed results. On the other hand, the
average relative error for the explicit inverse method remained
relatively low, around 5\,\%, signifying its superior ability to
handle the noise and produce more accurate reconstructions (see
Table~3). The values of regularization factors used in each case
from Fig.~19 and Fig.~20 are listed in Table~4.

% [page ends here — continues on page 13]

\end{document}
